Todos los casos son cliente-servidor sobre HTTP, pero cambian en la sincronía. En el SRI cada llamada es síncrona, aunque el trámite completo tiene dos pasos: primero se envía el comprobante a recepción y después se consulta su autorización con la clave de acceso \cite{sri_recepcion,sri_autorizacion}. Usa SOAP porque un solo WSDL sirve como contrato para cualquier sistema de facturación.

En nuestro PFC, Angular pregunta el estado cada 15 segundos. Es simple y funciona con el JWT, pero hace consultas aunque nada haya cambiado; con \emph{push} eso se evitaría. SGIRI muestra otra forma: un \texttt{POST} crea la medición en Globalping y luego un \texttt{GET} trae el resultado, así el cliente no queda bloqueado. DynamoDB usa HTTP y JSON, pero pone la operación en una cabecera y todo va por \texttt{POST} \cite{aws_dynamodb_api}: eso es RPC, no REST.
